\section{Binary System}

\subsection{The $X$--$P$--$T$ Diagram}
A \textcolor{blue}{binary system} has two components, $C=2$, so its composition is described by one
independent variable, $X=x_2$ with $x_1=1-X$. The phase diagram is therefore
three-dimensional in $(P,T,X)$. One can picture it as stacking the fixed-composition
$P$--$T$ diagrams and varying $X$ bit by bit. The phase rule reads
$F=4-\mathcal{P}$, and each unary structure is stretched by one dimension.

\begin{itemize}
    \item \textcolor{blue}{Tervariant volume} ($\mathcal{P}=1$, $F=3$). The bivariant regions of the
    unary diagram become volumes, where $P$, $T$ and $X$ can all be varied freely
    without changing phase. The cross section of such a volume changes with $X$.
    \item \textcolor{blue}{Bivariant surface} ($\mathcal{P}=2$, $F=2$). The univariant curves become
    surfaces. Two coexisting phases generally have different compositions, so each
    phase has its own bivariant surface, and the two surfaces are joined by tie
    lines. The region they enclose is the two-phase region, where the overall
    composition fixes the amount of each phase.
    \item \textcolor{blue}{Univariant curve} ($\mathcal{P}=3$, $F=1$). The invariant points become
    curves, along which only one variable can be varied freely. Three bivariant
    surfaces, one for each pair of phases, meet along each such curve.
    \item \textcolor{blue}{Invariant point} ($\mathcal{P}=4$, $F=0$). A new kind of invariant point
    appears, where none of $P$, $T$ or $X$ can be varied. Every subset of the four
    phases can coexist near it, so it is where $\binom{4}{3}=4$ univariant curves,
    $\binom{4}{2}=6$ two-phase (bivariant) surfaces and $\binom{4}{1}=4$
    single-phase (tervariant) volumes meet. Therefore, it is also called the \textcolor{blue}{quadruple point.}
\end{itemize}
